\section{Lie 군 관점의 이론적 정당화}

이 절의 목표는 ``PCA가 왜 좋은가''를 곧바로 말하는 것이 아니다. 먼저 모든 회전 기반 KV-cache 압축을 Lie 군 작용으로 다시 쓰고, 그 다음 정적 회전 선택 문제에 대한 보수적인 최적성 정리를 제시한다. 마지막으로 비등방적 비트 배분이 왜 자연스러운지 설명한다.

\subsection{회전은 Lie 군 위의 좌표계 선택이다}

\begin{definition}[회전 생성자]
\label{def:generator}
실수 직교 회전은 $SO(d)$의 원소 $R$로 나타내고, 그 Lie 대수는 반대칭 행렬의 공간
\[
\mathfrak{so}(d)=\{A \in \R^{d \times d}: A^\top=-A\}
\]
이다. $A \in \mathfrak{so}(d)$에 대해 $R(t)=\exp(tA)$는 $SO(d)$ 위의 일매개 부분군을 이룬다 \citep{hall2015lie}.
\end{definition}

\begin{proposition}[회전 기반 방법의 공통 형식]
\label{prop:common_form}
정적 직교 회전을 사용하는 모든 key 양자화 방법은 적절한 생성자 $A \in \mathfrak{so}(d)$와 시간 매개변수 $t$에 대해
\[
Q(k)=R\,\mathcal{Q}(R^\top k), \qquad R=\exp(tA),
\]
의 형태로 쓸 수 있다.
\end{proposition}

\begin{theorycomment}
\textbf{해설.} 이 명제는 새로운 수학 결과라기보다 문제를 정리하는 문장이다. random orthogonal rotation은 $A$를 무작위로 유도한 경우, identity는 $A=0$, PCA 회전은 데이터 공분산에 정렬된 기저가 유도하는 $R$, RoPE형 회전은 복소 2차원 평면의 블록별 생성자가 유도하는 unitary action의 실수 표현으로 읽을 수 있다.
\end{theorycomment}

이 관점이 중요한 이유는 ``회전이 좋다/나쁘다''는 문장을 더 세밀하게 바꾸기 때문이다. 좋은 회전이란 막연히 ``양자화 전에 한 번 돌려보는 것''이 아니라, \emph{문제에 맞는 생성자를 선택해 정보가 모이는 좌표계로 옮기는 것}이다. 따라서 이후 실험의 핵심도 ``PCA가 최종 PPL을 보장하는가''가 아니라 ``공분산 정렬 생성자가 무작위 생성자보다 더 유용한 surrogate를 주는가''가 된다.

\subsection{PCA는 공분산 정렬 정적 회전이다}

\begin{lemma}[상위 $k$축 분산 포획의 최적성]
\label{lem:kyfan}
대칭 양의 준정부호 행렬 $\Sigma \in \R^{d \times d}$와 정규직교 행렬 $U=[u_1,\dots,u_d]$를 생각하자. 임의의 $k<d$에 대해
\[
\sum_{i=1}^{k} u_i^\top \Sigma u_i
\]
를 최대화하는 $k$차원 부분공간은 $\Sigma$의 상위 $k$개 고유벡터가 생성하는 부분공간이다.
\end{lemma}

\begin{theorycomment}
\textbf{증명 개요.} 이는 Ky Fan 최대원리의 직접적 귀결이다 \citep{horn2012matrix}. 자세한 전개는 부록 A.1에 둔다.
\end{theorycomment}

\begin{theorem}[정적 회전 선택 문제에서의 PCA의 의미]
\label{thm:pca_optimal}
평균이 0인 랜덤 벡터 $x \in \R^d$의 공분산이 $\Sigma$라 하자. 모든 rank-$k$ 직교사영 $\Pi$ 가운데
\[
\E\lVert x-\Pi x\rVert_2^2
\]
를 최소화하는 사영은 $\Sigma$의 상위 $k$개 고유벡터에 대한 사영이다.
\end{theorem}

\begin{theorycomment}
\textbf{증명 개요.} 평균제곱 재구성 오차 최소화는 $\trace(\Pi\Sigma)$ 최대화와 동치이고, 이는 보조정리~\ref{lem:kyfan}로 바로 귀결된다. 자세한 증명은 부록 A.2에 둔다.
\end{theorycomment}

정리~\ref{thm:pca_optimal}가 말해 주는 것은 매우 제한적이지만 중요하다. 정적 선형 회전 가운데 \emph{에너지를 소수 축에 가장 잘 모으는 회전}을 찾는다면 PCA가 자연스럽다는 뜻이다. 이것이 곧바로 ``PCA-FOKVQ가 최종 PPL에서 항상 최고''를 뜻하지는 않는다. 그러나 적어도 정적 회전 문제의 첫 단계에서는 PCA를 임의 휴리스틱이 아니라 \emph{공분산 정렬 회전의 정답}으로 둘 수 있다.

\subsection{비트 배분은 회전 후의 별도 문제다}

\begin{assumption}[고율 축별 왜곡 근사]
\label{assump:high_rate}
회전된 좌표계에서 각 축이 약하게 분리되고, 축 $i$에 $b_i$ 비트를 할당했을 때 왜곡이
\[
D_i(b_i)\approx c_i \lambda_i 2^{-2b_i}
\]
로 근사된다고 가정한다. 여기서 $\lambda_i$는 해당 축의 분산이고 $c_i$는 분포 의존 상수다 \citep{cover2006elements}.
\end{assumption}

\begin{lemma}[water-filling 방향의 비트 배분]
\label{lem:waterfilling}
가정~\ref{assump:high_rate} 아래에서 총 비트 제약 $\sum_i b_i=B$ 하에
\[
D(\mathbf{b})=\sum_i c_i\lambda_i 2^{-2b_i}
\]
를 최소화하는 active set 내부 해는
\[
b_i=\alpha+\frac{1}{2}\log_2(c_i\lambda_i)
\]
형태를 갖는다. 여기서 $\alpha$는 예산 제약을 맞추는 상수다.
\end{lemma}

\begin{theorycomment}
\textbf{증명 개요.} 라그랑지안 최적화로 각 active axis의 한계 왜곡이 같아지며, 자세한 유도는 부록 A.3에 둔다.
\end{theorycomment}

이 보조정리는 FOKVQ의 gamma 스케줄이 문자 그대로 최적이라고 말해 주지 않는다. 대신 \emph{분산이 큰 축에 더 많은 비트를 주는 방향성 자체가 왜 자연스러운가}를 설명한다. 다시 말해, Lie 군 해석은 회전 선택을 정당화하고, water-filling 근사는 그 회전된 좌표에서의 비트 배분 방향을 정당화한다. 현재 구현은 이 두 층을 거친 휴리스틱에 가깝다.

\begin{remark}
이 절의 핵심 분리는 끝까지 유지되어야 한다. \textbf{회전의 질}과 \textbf{회전 후 양자화기의 질}은 같은 문제가 아니다. 현재 실험에서 PCA가 random/identity보다 나은 것은 첫 번째를 지지한다. 반면 KIVI류 baseline을 넘지 못하는 것은 두 번째가 여전히 병목일 수 있음을 시사한다.
\end{remark}
